\chapter{Pattern-Match-Oriented Programming Exercises}

To help you experience pattern-match-oriented programming, we have prepared problems in quiz format.

\section*{The \texttt{member} Function}

Fill in the blank below to complete the \lstinline{member} function.

\begin{lstlisting}[language=egison]
def member x xs :=
  match xs as list eq with
  | `\fcolorbox{black}[gray]{1.0}{\vphantom{0}\hspace{5em}1\vphantom{0}\hspace{5em}}` -> True
  | _ -> False

member 1 [1,2,3,4]
-- True

member 5 [1,2,3,4]
-- False
\end{lstlisting}

\section*{The \texttt{concat} Function}

Fill in the blanks below to complete the \lstinline{concat} function.

\begin{lstlisting}[language=egison]
def concat xss :=
  matchAllDFS xss as `\fcolorbox{black}[gray]{1.0}{\vphantom{0}\hspace{5em}1\vphantom{0}\hspace{5em}}` with
  | `\fcolorbox{black}[gray]{1.0}{\vphantom{0}\hspace{5em}2\vphantom{0}\hspace{5em}}` -> `\fcolorbox{black}[gray]{1.0}{\vphantom{0}\hspace{5em}3\vphantom{0}\hspace{5em}}`

concat [[1,2],[3],[4,5]]
-- [1,2,3,4,5]
\end{lstlisting}

\section*{Quadruplet Primes}

Write a program that enumerates quadruplet primes.
Quadruplet primes are quadruples of primes of the form $(p,p + 2,p + 6,p + 8)$.

\begin{lstlisting}[language=egison]
take 4 (matchAll primes as list integer with
        | `\fcolorbox{black}[gray]{1.0}{\vphantom{0}\hspace{5em}1\vphantom{0}\hspace{5em}}` -> (p, p + 2, p + 6, p + 8))
-- [(5, 7, 11, 13), (11, 13, 17, 19), (101, 103, 107, 109), (191, 193, 197, 199)]
\end{lstlisting}

\section*{Prime Pairs with Difference $6$}

Fill in the blank below to extract prime pairs $(p, p+6)$.

\begin{lstlisting}[language=egison]
take 6 (matchAll primes as list integer with
        | `\fcolorbox{black}[gray]{1.0}{\vphantom{0}\hspace{5em}1\vphantom{0}\hspace{5em}}` -> (p, p + 6))
-- [(5, 11), (7, 13), (11, 17), (13, 19), (17, 23), (23, 29)]
\end{lstlisting}

\section*{The \texttt{intersect} Function}

Fill in the blanks below to complete the \lstinline{intersect} function, which extracts the intersection of two collections.

\begin{lstlisting}[language=egison]
def intersect xs ys :=
  matchAllDFS (xs, ys) as `\fcolorbox{black}[gray]{1.0}{\vphantom{0}\hspace{5em}1\vphantom{0}\hspace{5em}}` with
  | `\fcolorbox{black}[gray]{1.0}{\vphantom{0}\hspace{5em}2\vphantom{0}\hspace{5em}}` -> `\fcolorbox{black}[gray]{1.0}{\vphantom{0}\hspace{5em}3\vphantom{0}\hspace{5em}}`

intersect [1,2,3,4] [2,4,5]
-- [2,4]
\end{lstlisting}

\section*{The \texttt{difference} Function}

Fill in the blanks below to complete the \lstinline{difference} function, which extracts only the elements of the first collection that are not contained in the second collection.

\begin{lstlisting}[language=egison]
def difference xs ys :=
  matchAllDFS (xs, ys) as `\fcolorbox{black}[gray]{1.0}{\vphantom{0}\hspace{5em}1\vphantom{0}\hspace{5em}}` with
  | `\fcolorbox{black}[gray]{1.0}{\vphantom{0}\hspace{5em}2\vphantom{0}\hspace{5em}}` -> `\fcolorbox{black}[gray]{1.0}{\vphantom{0}\hspace{5em}3\vphantom{0}\hspace{5em}}`

difference [1,2,3,4] [2,4,5]
-- [1,3]
\end{lstlisting}

\section*{The \texttt{doubles} Function}

Complete the \lstinline{doubles} function, which extracts only the elements that appear exactly $2$ times in a collection.

\begin{lstlisting}[language=egison]
def doubles xs :=
  matchAllDFS xs as `\fcolorbox{black}[gray]{1.0}{\vphantom{0}\hspace{5em}1\vphantom{0}\hspace{5em}}` with
  | `\fcolorbox{black}[gray]{1.0}{\vphantom{0}\hspace{5em}2\vphantom{0}\hspace{5em}}` -> `\fcolorbox{black}[gray]{1.0}{\vphantom{0}\hspace{5em}3\vphantom{0}\hspace{5em}}`

doubles [1,2,3,2,1,1,4,5,3]
-- [2,3]
\end{lstlisting}

\section*{The \texttt{tails} Function}

Define the \lstinline{tails} function using pattern matching.
Try to find two solutions: one using loop patterns and one without.

\begin{lstlisting}[language=egison]
def tails xs :=
  matchAll xs as list something with
  | `\fcolorbox{black}[gray]{1.0}{\vphantom{0}\hspace{5em}1\vphantom{0}\hspace{5em}}`
  -> ts

tails [1,2,3,4]
-- [[1, 2, 3, 4], [2, 3, 4], [3, 4], [4], []]
\end{lstlisting}

\section*{Multiset Equality Check}

Define a two-argument predicate \lstinline{equalMultiset} that checks multiset equality using pattern matching.

\begin{lstlisting}[language=egison]
def equalMultiset xs ys :=
  match (xs, ys) as (list eq, multiset eq) with
  | ([], [])
  -> True
  | `\fcolorbox{black}[gray]{1.0}{\vphantom{0}\hspace{5em}1\vphantom{0}\hspace{5em}}`
  -> equalMultiset xs' ys'
  | _
  -> False

equalMultiset [1,1,2] [2,1,1]
-- True

equalMultiset [1,1,2] [1,2,2]
-- False
\end{lstlisting}

Try writing it without recursion using loop patterns.

\begin{lstlisting}[language=egison]
def equalMultiset xs ys :=
  match (xs, ys) as (list eq, multiset eq) with
  | `\fcolorbox{black}[gray]{1.0}{\vphantom{0}\hspace{5em}2\vphantom{0}\hspace{5em}}`
  -> True
  | _
  -> False
\end{lstlisting}

\section*{Poker Hand Evaluation Considering Jokers}

Without modifying the implementation of the poker hand evaluation pattern matching (\lstinline{poker} function) from Chapter~\ref{pmo}, Section~\ref{pmo-poker} at all, modify only the definition of the \lstinline{card} matcher to enable poker hand evaluation that accounts for jokers.
\begin{lstlisting}[language=egison]
def card := matcher
  | card $ $ as (suit, mod 13) with 
    | Card $s $n -> [(s, n)]
    | Joker -> `\fcolorbox{black}[gray]{1.0}{\vphantom{0}\hspace{5em}1\vphantom{0}\hspace{5em}}`
  | $ as something with
    | $tgt -> [tgt]

poker [Card Spade 5, Card Spade 6, Joker, Card Spade 8, Card Spade 9]
-- "Straight flush"
poker [Card Spade 5, Card Diamond 5, Joker, Card Club 5, Card Heart 7]
-- "Four of a kind"
\end{lstlisting}
